\begin{statementbox}
	\textbf{\textcolor{CStatement}{条件}}
	\begin{description}[style=nextline, leftmargin=2.8em, labelsep=0.5em, itemsep=0.35em]
		\item[\textbf{\textcolor{CStatement}{条件1（小量近似）：}}] 自变量 $x$ 满足 $x\to0$，实际估算时通常要求 $|x|$ 较小。
		\item[\textbf{\textcolor{CStatement}{条件2（弧度制）：}}] 涉及 $\sin x,\cos x,\tan x$ 时，$x$ 必须用弧度制表示。
		\item[\textbf{\textcolor{CStatement}{条件3（定义域与光滑性）：}}] 使用 $\ln(1+x)$ 时需保证 $1+x>0$；使用 $\tan x$ 时需保证 $\cos x\ne0$；一般函数展开要求在 $x=0$ 附近足够光滑。
	\end{description}

	\textbf{\textcolor{CStatement}{结论}}
	\begin{description}[style=nextline, leftmargin=2.8em, labelsep=0.5em, itemsep=0.45em]
		\item[\textbf{\textcolor{CStatement}{结论一（指数近似）：}}]
		      \textbf{\textcolor{CStatement}{直观描述：}} $e^x$ 在零点附近的二次近似。
		      \[
			      e^x = 1 + x + \frac{x^2}{2} + O(x^3)
		      \]
		      因此常用近似为
		      \[
			      e^x \approx 1+x+\frac{x^2}{2}.
		      \]

		\item[\textbf{\textcolor{CStatement}{结论二（对数近似）：}}]
		      \textbf{\textcolor{CStatement}{直观描述：}} $\ln(1+x)$ 在零点附近的三次近似。
		      \[
			      \ln(1+x)=x-\frac{x^2}{2}+\frac{x^3}{3}+O(x^4),\qquad x>-1.
		      \]
		      因此常用近似为
		      \[
			      \ln(1+x)\approx x-\frac{x^2}{2}+\frac{x^3}{3}.
		      \]

		\item[\textbf{\textcolor{CStatement}{结论三（正弦近似）：}}]
		      \textbf{\textcolor{CStatement}{直观描述：}} $\sin x$ 在零点附近的五次近似，其中 $x$ 为弧度制。
		      \[
			      \sin x=x-\frac{x^3}{6}+\frac{x^5}{120}+O(x^7)
		      \]
		      因此常用近似为
		      \[
			      \sin x\approx x-\frac{x^3}{6}+\frac{x^5}{120}.
		      \]

		\item[\textbf{\textcolor{CStatement}{结论四（余弦近似）：}}]
		      \textbf{\textcolor{CStatement}{直观描述：}} $\cos x$ 在零点附近的四次近似，其中 $x$ 为弧度制。
		      \[
			      \cos x=1-\frac{x^2}{2}+\frac{x^4}{24}+O(x^6)
		      \]
		      因此常用近似为
		      \[
			      \cos x\approx 1-\frac{x^2}{2}+\frac{x^4}{24}.
		      \]

		\item[\textbf{\textcolor{CStatement}{结论五（正切近似）：}}]
		      \textbf{\textcolor{CStatement}{直观描述：}} $\tan x$ 在零点附近的三次近似，其中 $x$ 为弧度制。
		      \[
			      \tan x=x+\frac{x^3}{3}+O(x^5)
		      \]
		      因此常用近似为
		      \[
			      \tan x\approx x+\frac{x^3}{3}.
		      \]

		\item[\textbf{\textcolor{CStatement}{推广结论（泰勒通式）：}}]
		      \textbf{\textcolor{CStatement}{直观描述：}} 足够光滑函数在零点附近的多项式逼近形式。
		      \[
			      f(x)=f(0)+f'(0)x+\frac{f''(0)}{2!}x^2+\cdots+\frac{f^{(n)}(0)}{n!}x^n+o(x^n),\qquad x\to0.
		      \]
	\end{description}
\end{statementbox}